% ---------------------------------------------------------------------------
% 4 The plaintext reference
% Source: docs/tasks.md, docs/eval_approach.md, docs/data_provenance.md
%         numbers: evidence/measurements.yaml plaintext_baseline
% ---------------------------------------------------------------------------

\section{Plaintext validation}
\label{sec:baseline}

Before encrypted evaluation, we reproduced DNAGPT on genomic-signal recognition, promoter and
splice-site classification, and mRNA abundance regression. We also recorded the plaintext outputs
required for subsequent numerical comparison with CKKS execution.

The upstream DNAGPT implementation is used without changing its forward computation. Published or
reconstructed evaluation splits are used for the selected tasks. Genomic-signal recognition is
evaluated on the full balanced source corpus, including examples used to train its head, so this
result measures model and pipeline reproduction rather than held-out generalization.

Numerical validation uses two independent float64 implementations. A PyTorch execution and a
separate NumPy implementation are compared after every transformer block and the task head, with a
$2\times10^{-5}$ agreement criterion. After this check passes, the NumPy implementation provides
the plaintext reference for encrypted evaluation. The original float32 execution is retained only
to quantify dtype-related drift.

\begin{table*}[!tb]\centering\small
\caption{Plaintext DNAGPT validation on three genomic task families. Signal recognition and
abundance regression use the released fine-tuned heads; the understanding-benchmark rows use a
locally fine-tuned head because DNAGPT provides no corresponding head.}
\label{tab:baseline}
\begin{tabularx}{\linewidth}{Y l l l l}
\toprule
Task & Metric & This work & Reference & Source \\
\midrule
Polyadenylation-signal recognition$^{\ddagger}$ & accuracy & 0.9124 & 0.9151 & \citet{dnagpt2023} \\
 & F1 & 0.916 & --- & \\
Core promoter detection & MCC & 0.680 & 0.69 & \citet{dnabert2} \\
300\,bp promoter detection & MCC & 0.897 & 0.87 & \citet{dnabert2} \\
Splice-site classification & MCC & 0.831 & 0.85 & \citet{dnabert2} \\
mRNA abundance regression$^{\dagger\S}$ & $r^2$ & 0.562 & 0.62 & \citet{dnagpt2023} \\
 & Pearson $r$ & 0.753 & --- & \\
\bottomrule
\end{tabularx}
\vspace{2pt}
{\footnotesize $^{\ddagger}$The reference is DNAGPT's own reported accuracy for this model class,
measured on a held-out quarter of the set; this work evaluates all $22{,}604$ examples. The
comparison is same-model but not same-split.\\
$^{\dagger}$Model-fidelity evidence only. This task is outside the encrypted
scope; see Section~\ref{sec:limitations}.\\
$^{\S}$The released regression head evaluated here is fine-tuned from the human-pretrained
$0.1$B variant, while the reference is reported for the mammal-pretrained variant, so the
backbones differ.}
\end{table*}

Across the selected tasks, the reproduced results retain the expected task signal and provide the
model-validation baseline used before encrypted evaluation. The GUE tasks use locally trained heads
because DNAGPT provides no corresponding fine-tuned head, and the recovered abundance split is used
only as model-fidelity evidence.

\subsection{Benchmark prompt length}
Genomic-signal recognition uses $600$-base-pair inputs. DNAGPT tokenizes each input into $100$
$6$-mer sequence tokens plus three task-format tokens, yielding a $103$-token prompt. This is the
longest prompt among the encrypted-scope classification tasks and is therefore used as the
maximum-length encrypted benchmark in Sections~\ref{sec:protocol} and~\ref{sec:results}.
